\section{Anexos}
\label{sec:anexos}

Esta sección reúne el material de apoyo exigido por la guía de la Entrega Final que no encaja como
prosa dentro de los capítulos anteriores. Donde el material solicitado no existe todavía de forma real,
se declara la brecha explícitamente en vez de fabricarlo --- la misma política de integridad aplicada en
el resto de este informe (Sección~\ref{sec:discusion}).

\subsection{Anexo A --- Cadena de búsqueda de la revisión de trabajos relacionados}
\label{anexo:busqueda}

\textbf{Brecha declarada.} La búsqueda dirigida de la Sección~\ref{sec:trabajos-relacionados} reporta el
número de consultas ejecutadas (24) y el número de resultados aproximado (\textasciitilde200), pero
\textbf{no se registraron las cadenas de búsqueda booleanas exactas ni la fecha de cada consulta} en su
momento --- así lo declara \texttt{docs/checklists/PRISMA-2020.md} en el criterio ``Estrategia de
búsqueda completa y reproducible'': \textit{``se reporta el número de consultas (24) pero no las
cadenas de búsqueda exactas ni las fechas de ejecución --- no es reproducible palabra por palabra''}. No
se reconstruye aquí una cadena de búsqueda aproximada a partir de memoria, porque presentarla como si
fuera la consulta real ejecutada sería exactamente el tipo de dato fabricado que este informe corrige en
otros lugares (Secciones~\ref{sec:evaluacion} y \ref{sec:discusion}). Quedan como palabras clave
efectivamente usadas, sin cadena booleana ni fecha verificable: \textit{software requirements
traceability}, \textit{JWT authentication Spring Boot}, \textit{stored procedures vs ORM performance},
\textit{System Usability Scale}, \textit{OWASP dependency scanning}, \textit{CI/CD reproducibility},
entre otras relacionadas con cada uno de los 13 trabajos comparados en la Tabla de la
Sección~\ref{sec:trabajos-relacionados}.

\subsection{Anexo B --- Protocolo del instrumento SUS (System Usability Scale)}
\label{anexo:sus}

El instrumento sigue la forma estándar de Brooke~\cite{brooke1996}: diez enunciados de alternancia de
polaridad (impares en sentido positivo, pares en sentido negativo), cada uno respondido en una escala
Likert de 1 (totalmente en desacuerdo) a 5 (totalmente de acuerdo).

\begin{enumerate}[nosep]
    \item Creo que me gustaría utilizar este sistema con frecuencia.
    \item Encontré el sistema innecesariamente complejo.
    \item Pensé que el sistema era fácil de usar.
    \item Creo que necesitaría el apoyo de un técnico para poder utilizar este sistema.
    \item Encontré que las diversas funciones del sistema estaban bien integradas.
    \item Pensé que había demasiada inconsistencia en el sistema.
    \item Imagino que la mayoría de las personas aprenderían a usar el sistema muy rápidamente.
    \item Encontré el sistema muy pesado/incómodo de usar.
    \item Me sentí muy confiado/seguro al usar el sistema.
    \item Necesité aprender muchas cosas antes de poder empezar a usar el sistema.
\end{enumerate}

\textbf{Cálculo de la puntuación:} para los ítems impares (1,3,5,7,9) se resta 1 a la respuesta; para los
pares (2,4,6,8,10) se resta la respuesta de 5. La suma de las diez contribuciones se multiplica por 2.5,
dando un puntaje individual entre 0 y 100 (no es un porcentaje). El script
\texttt{scripts/sus-analysis.ipynb} implementa exactamente esta fórmula y trae tres casos de autotest que
pasan.

\textbf{Estado real (Sección~\ref{sec:evaluacion}):} el instrumento está listo pero \textbf{no se ha
aplicado a participantes reales} ($N=0$). El protocolo de consentimiento informado individual para cuando
se aplique está redactado como plantilla en \texttt{docs/etica/consentimientos/plantilla-consentimiento.md}
(explica qué se recoge, para qué, cuánto tiempo se conserva, quién lo verá y cómo se retira el
consentimiento), pero no existe todavía ninguna aprobación fechada previa a una recogida real, ni
participantes reclutados. Una versión anterior de este proyecto publicó un puntaje fabricado de
91.25/100 con diez evaluadores inexistentes; esa cifra fue retirada explícitamente y se documenta en
\texttt{docs/mediciones/sus/SUS-RESULTS.md} y en la Sección~\ref{sec:evaluacion} de este informe.

\subsection{Anexo C --- Corridas de integración continua}
\label{anexo:ci}

\textbf{Brecha declarada.} La guía pide capturas de tres corridas verdes de CI. El workflow
\texttt{.github/workflows/ci.yml} existe, corre en cada \textit{push} y ejecuta build, tests, JaCoCo y
análisis estático de seguridad del backend contra Postgres/Redis reales (no simulados). Sin embargo, no
se generaron ni versionaron capturas de pantalla de corridas reales del panel de GitHub Actions durante
la elaboración de este informe, y no fue posible capturarlas al cierre de esta corrección por una falla
temporal de la herramienta de navegador disponible en este entorno de trabajo. En vez de sustituir esta
evidencia con una captura simulada o de otro proyecto, se declara la brecha: falta que el equipo entre a
\url{https://github.com/carla22072004/PFC-Presustentaciones-2026/actions}, localice tres corridas
recientes con estado \textit{success} en la rama que se va a defender, y capture pantalla de cada una
antes del examen final.

\subsection{Anexo D --- Checklist de estándares empíricos de Ralph et al. (2021)}
\label{anexo:ralph}

Además de \texttt{docs/checklists/RALPH-2021-GENERAL.md} (Estándar General, resumido en la
Tabla~\ref{tab:ralph-general}), el repositorio aplica dos checklists adicionales de la misma familia
--- \texttt{RALPH-2021-BENCHMARKING.md} (para las mediciones de rendimiento con k6) y
\texttt{RALPH-2021-ENGINEERING-RESEARCH.md} (para el proceso de desarrollo como tal) --- referenciados
desde la Sección~\ref{sec:marco-teorico} pero no reproducidos aquí en extenso por brevedad.

\begin{table}[H]
\centering
\small
\begin{tabular}{p{7.5cm}p{1.3cm}p{5cm}}
\toprule
\textbf{Criterio (Estándar General)} & \textbf{Cumple} & \textbf{Evidencia resumida} \\
\midrule
Afirmaciones empíricas basadas en datos, no en opinión & Sí & Toda cifra citada tiene archivo crudo de origen (\texttt{DATA-PROVENANCE.md}) \\
Método de recolección reproducible (comandos documentados) & Sí & Comandos exactos en \texttt{k6/README.md}, \texttt{LIGHTHOUSE-REPORT.md}, \texttt{OWASP-AUDIT.md} \\
Contexto/entorno de medición declarado & Sí & \texttt{docs/entorno/versions.txt} + notas metodológicas explícitas \\
Se reportan resultados negativos, no solo favorables & Sí & Endpoint inexistente en k6, 233 hallazgos find-sec-bugs, 14 vulnerabilidades npm sin corregir \\
Correlación vs. causalidad distinguidas & Parcial & Cache fría/caliente sí aísla variable; runs 1--2 de k6 no lo hacen, y el reporte lo aclara \\
Amenazas a la validez discutidas sistemáticamente & No & Sin sección dedicada por reporte individual (sí existe una transversal en Sección~\ref{sec:amenazas-validez} de este informe) \\
Tamaño de muestra justificado & Parcial & $n=30$ en caché sigue convención CLT sin justificación explícita en texto; 5 corridas k6 siguen el mínimo de la guía, no un cálculo de poder \\
Artefactos disponibles para verificación independiente & Sí & \texttt{DATA-PROVENANCE.md} enlaza cada archivo crudo referenciado \\
\bottomrule
\end{tabular}
\caption{Checklist de estándares empíricos generales de Ralph et al.~\cite{ralph2021} --- 5/8 cumplidos, 2 parciales, 1 no cumplido.}
\label{tab:ralph-general}
\end{table}

\subsection{Anexo E --- Matriz de trazabilidad completa}
\label{anexo:matriz}

La Tabla~\ref{tab:matriz-completa} reproduce íntegramente \texttt{docs/trazabilidad/matriz.csv} (15
filas, 9 columnas), validada por \texttt{scripts/validate-traceability.sh} y por un parseo estricto con
\texttt{csv.reader} de Python que confirma que las 15 filas tienen exactamente 9 campos cada una ---
cinco de ellas (RF-05, RF-07, RF-08, RF-09, RF-10) tenían comas sin escapar en la columna de evidencia
empírica que rompían el conteo de columnas; se corrigieron entre comillas en la corrección del 2026-09-05
sin alterar el contenido de la evidencia. La columna ``Módulo'' se abrevia en algunas filas por espacio;
el CSV original tiene el nombre completo.

\input{secciones/matriz-completa-generada}
